% =====================================================================
%  Chapter 27 — The Certainty-Equivalent Tilt
% =====================================================================

\chapter{The Certainty-Equivalent Tilt}\idx{Tilt}
\label{ch:tilt}

Chapter 26 derived a sizing\idx{Sizing} rule for assets. This chapter applies it to
the book's actual problem --- tilting inside a sleeve, without touching
the sleeve --- and explains the mechanics that turn a set of position
sizes into a portfolio the investor can hold.

\section{Why the tilt is bounded and renormalised}

The certainty-equivalent solve of Chapter 26 produces, for each sleeve, a
vector of weights. Left alone, those weights would replace the sleeve's
allocation entirely, which is not what the book wants: the sleeve's
\emph{total} is set by policy, and the model is permitted to change only
its composition.

The book therefore does three things in sequence.

First, it takes the certainty-equivalent weights for the sleeve's live
members and clips the deviations from the base weights to a bound:
\begin{equation}
d_i = \operatorname{clip}\!\left(
  \frac{w_i^{CE}}{b_i} - 1,\ -\frac{1}{2},\ \frac{1}{2}\right),
\end{equation}
where \(b_i\) is the base weight of member \(i\) inside the sleeve.

Second, it applies the tilt:
\begin{equation}
\tilde{b}_i = b_i\left(1 + \kappa\, d_i\right),
\end{equation}
where \(\kappa\) is the tilt scale. The book uses \(\kappa = 0.25\), so the
tilt moves a sleeve's composition by at most an eighth in either
direction.

Third, it renormalises the sleeve back to its policy weight:
\begin{equation}
b_i \leftarrow \tilde{b}_i\,\frac{\sum_j b_j}{\sum_j \tilde{b}_j},
\end{equation}
so that the sleeve's total is exactly what the policy says.

\begin{takeaway}[What the bound\idx{Bound} buys]
The tilt can change which funds inside a sleeve are held and can move the
sleeve's total within a bound, and it can do neither by much. A reader
who disagrees with the model can undo its entire contribution by
reverting to the base weights, and the difference will be at most a
fraction of a per cent of the portfolio per month.
\end{takeaway}

\section{Why the clipping is a second layer of protection}

The clip at \(\pm 1/2\) on the deviation is not a numerical convenience.
It bounds the influence of any single asset's forecast on the portfolio.

Without it, a single fund whose posterior mean happened to be extreme
would dominate its sleeve's allocation, and the tilt would stop being a
tilt and become a concentrated bet on one forecast. With it, the maximum
sleeve-relative move per member is 50 per cent before the tilt scale is
applied, and 12.5 per cent after.

The second-order effect is more important than the first, and it is
favourable. The clip is a *non-linearity* that the certainty-equivalent
rule did not ask for, and it makes the tilt's response to an extreme
forecast flatter than the linear rule would be. The practical consequence
is that the tilt's contribution to the portfolio's return distribution has
bounded tails, which is what makes it safe to add to a portfolio whose
main job is to protect principal.

\section{What the tilt does in practice}

The book's system produces an average cash weight and an average turnover,
both reported in Part VI, and they are the direct evidence of how much the
tilt does.

The qualitative picture is stable across the whole record and across
almost every configuration in the trial family. The tilt is small: the
average sleeve-relative deviation is a fraction of a percentage point.
The turnover is modest: a few per cent a month. The mean cash weight is
small but non-zero, and it rises in the periods when the posterior over
the edge is least favourable.

This is the expected behaviour of a system whose forecast is a posterior
with a large variance, and Chapter 20's arithmetic predicts it exactly. A
forecast that is a fraction of a standard deviation from zero produces a
certainty-equivalent weight that is nearly zero.

The response of the certainty-equivalent weight to the signal-to-noise ratio of the forecast is shown in \figref{fig:kelly-shrinkage}.

\section{The cost of the tilt}

Two costs, and both are small.

The \emph{spread cost} is charged on every one-way trade, at half the
quoted spread. The tilt is the source of essentially all of the system's
turnover, so the cost of the tilt is the cost of the whole system.

The \emph{tax cost} is subtler and larger. The tilt causes the portfolio
to realise gains and losses, and realised gains are taxed at the
preferential rate while realised losses offset ordinary income only up to
\$3,000 a year. A portfolio that tilts has a systematically higher
turnover of realised gains, and the tax bill is the price of the tilt.

The book's taxable account exists to measure exactly this, and Part VI
reports the gap between the two accounts. It is the clearest single
illustration of the point of Chapter 3: a strategy that does not work has
costs, and the costs are paid whether or not the strategy pays.

\begin{keytakeaways}
\item The tilt is clipped, scaled and renormalised, so that the model can
change a sleeve's composition but not its total, and can move the total
only within a bound.
\item The clip is a second, non-linear layer of protection that bounds the
tilt's contribution to the tail of the portfolio's return distribution.
\item The tilt's actual magnitude is a fraction of a percentage point per
sleeve, which is what Chapter 20's arithmetic predicts.
\item The tilt's costs are real: it is the source of the turnover and the
tax bill, and it pays them whether or not it pays the investor.
\end{keytakeaways}
